\begin{textsample}{POS Dim 1 – 2005 – Score 27.00 – 2005/t000476.txt}  \label{ex:f1_pos_006}
Well, it ' s great to be here.
We ' ve heard a lot about \textbf{the} promise of technology, and \textbf{the} peril.
I ' ve been quite interested in both.
If we could convert 0. 03 percent of \textbf{the} sunlight that falls on \textbf{the} earth into energy, we could meet all of our projected needs for 2030.
We can ' t do that today because solar panels are heavy, expensive and very inefficient.
There are nano-engineered designs, which at least have been analyzed theoretically, that show \textbf{the} potential to be very lightweight, very inexpensive, very efficient, and we ' d be able to actually provide all of our energy needs in this \textbf{renewable} way.
Nano-engineered fuel cells could provide \textbf{the} energy where it ' s needed.
That ' s a key trend, which is decentralization, moving from centralized nuclear power plants and liquid natural gas tankers to decentralized resources that are environmentally more \textbf{friendly}, a lot more efficient and capable and safe from disruption.
Bono spoke very eloquently, that we have \textbf{the} tools, for \textbf{the} first time, to address age-old problems of disease and poverty.
Most regions of \textbf{the} world are moving in that direction.
In 1990, in East Asia and \textbf{the} Pacific region, there were 500 million people living in poverty — that number now is under 200 million. \textbf{The} World Bank projects by 2011, it will be under 20 million, which is a reduction of 95 percent.
I did \textbf{enjoy} Bono ' s comment linking Haight-Ashbury to Silicon Valley.
Being from \textbf{the} Massachusetts high-tech community myself, I ' d point out that we were hippies also in \textbf{the} 1960s, although we hung around Harvard Square.
But we do have \textbf{the} potential to overcome disease and poverty, and I ' m going to talk about those issues, if we have \textbf{the} will.
Kevin Kelly talked about \textbf{the} acceleration of technology.
That ' s been a strong interest of mine, and a \textbf{theme} that I ' ve developed for some 30 years.
I realized that my technologies had to make sense when I finished a project.
That invariably, \textbf{the} world was a different place when I would introduce a technology.
And, I noticed that most inventions fail, not because \textbf{the} R\&D department can ' t get it to work — if you look at most business plans, they will actually succeed if given \textbf{the} opportunity to build what they say they ' re going to build — and 90 percent of those projects or more will fail, because \textbf{the} timing is wrong — not all \textbf{the} enabling factors will be in place when they ' re needed.
So I began to be an ardent student of technology trends, and track where technology would be at different points in time, and began to build \textbf{the} mathematical models of that.
It ' s kind of taken on a life of its own.
I ' ve got a group of 10 people that work with me to gather data on key measures of technology in many different areas, and we build models.
And you ' ll hear people say, well, we can ' t predict \textbf{the} future.
And if you ask me, will \textbf{the} price of Google be higher or lower than it is today three years from now, that ' s very hard to say.
Will WiMax CDMA G3 be \textbf{the} wireless standard three years from now?
That ' s hard to say.
But if you ask me, what will it cost for one MIPS of computing in 2010, or \textbf{the} cost to sequence a base pair of DNA in 2012, or \textbf{the} cost of sending a megabyte of data wirelessly in 2014, it turns out that those are very predictable.
There are remarkably smooth exponential curves that govern price performance, capacity, bandwidth.
And I ' m going to show you a small sample of this, but there ' s really a theoretical reason why technology develops in an exponential fashion.
And a lot of people, when they think about \textbf{the} future, think about it linearly.
They think they ' re going to continue to develop a problem or address a problem using today ' s tools, at today ' s pace of progress, and fail to take into consideration this exponential growth. \textbf{The} \textbf{Genome} Project was a controversial project in 1990.
We had our best Ph.
D. students, our most advanced equipment around \textbf{the} world, we got 1 / 10, 000th of \textbf{the} project done, so how ' re we going to get this done in 15 years?
And 10 years into \textbf{the} project, \textbf{the} skeptics were still going strong — says,"You ' re two- thirds through this project, and you ' ve managed to only sequence a very tiny percentage of \textbf{the} whole \textbf{genome}."But it ' s \textbf{the} nature of exponential growth that once it reaches \textbf{the} knee of \textbf{the} curve, it explodes.
Most of \textbf{the} project was done in \textbf{the} last few years of \textbf{the} project.
It took us 15 years to sequence HIV — we sequenced SARS in 31 days.
So we are gaining \textbf{the} potential to overcome these problems.
I ' m going to show you just a few examples of how pervasive this phenomena is. \textbf{The} actual paradigm-shift rate, \textbf{the} rate of adopting new ideas, is doubling every decade, according to our models.
These are all logarithmic graphs, so as you go up \textbf{the} levels it represents, generally multiplying by factor of 10 or 100.
It took us half a century to adopt \textbf{the} telephone, \textbf{the} first virtual- reality technology.
Cell phones were adopted in about eight years.
If you put different communication technologies on this logarithmic graph, television, radio, telephone were adopted in decades.
Recent technologies — like \textbf{the} PC, \textbf{the} \textbf{web}, cell phones — were under a decade.
Now this is an interesting chart, and this really gets at \textbf{the} fundamental reason why an \textbf{evolutionary} process — and both \textbf{biology} and technology are \textbf{evolutionary} processes — accelerate.
They work through interaction — they create a capability, and then it uses that capability to bring on \textbf{the} next stage.
So \textbf{the} first step in biological \textbf{evolution}, \textbf{the} \textbf{evolution} of DNA — actually it was RNA came first — took billions of years, but then \textbf{evolution} used that information-processing backbone to bring on \textbf{the} next stage.
So \textbf{the} Cambrian \textbf{Explosion}, when all \textbf{the} body plans of \textbf{the} animals were \textbf{evolved}, took only 10 million years.
It was 200 times faster.
And then \textbf{evolution} used those body plans to \textbf{evolve} higher cognitive functions, and biological \textbf{evolution} kept accelerating.
It ' s an inherent nature of an \textbf{evolutionary} process.
So Homo sapiens, \textbf{the} first technology-creating species, \textbf{the} species that combined a cognitive function with an opposable appendage — and by \textbf{the} way, chimpanzees don ' t really have a very good opposable thumb — so we could actually \textbf{manipulate} our environment with a power grip and fine motor coordination, and use our mental models to actually change \textbf{the} world and bring on technology.
But anyway, \textbf{the} \textbf{evolution} of our species took hundreds of thousands of years, and then working through interaction, \textbf{evolution} used, essentially, \textbf{the} technology-creating species to bring on \textbf{the} next stage, which were \textbf{the} first steps in technological \textbf{evolution}.
And \textbf{the} first step took tens of thousands of years — stone tools, fire, \textbf{the} wheel — kept accelerating.
We always used then \textbf{the} latest generation of technology to create \textbf{the} next generation.
Printing press took a century to be adopted ; \textbf{the} first \textbf{computers} were designed pen-on-paper — now we use \textbf{computers}.
And we ' ve had a continual acceleration of this process.
Now by \textbf{the} way, if you look at this on a linear graph, it looks like everything has just happened, but some observer says,"Well, Kurzweil just put points on this graph that fall on that straight line."So, I took 15 different lists from key thinkers, like \textbf{the} Encyclopedia Britannica, \textbf{the} Museum of Natural History, Carl Sagan ' s Cosmic Calendar on \textbf{the} same — and these people were not trying to make my point ; these were just lists in reference works, and I think that ' s what they thought \textbf{the} key events were in biological \textbf{evolution} and technological \textbf{evolution}.
And again, it forms \textbf{the} same straight line.
You have a little bit of thickening in \textbf{the} line because people do have disagreements, what \textbf{the} key points are, there ' s differences of opinion when \textbf{agriculture} started, or how long \textbf{the} Cambrian \textbf{Explosion} took.
But you see a very clear trend.
There ' s a basic, profound acceleration of this \textbf{evolutionary} process.
Information technologies double their capacity, price performance, bandwidth, every year.
And that ' s a very profound \textbf{explosion} of exponential growth.
A personal experience, when I was at MIT — \textbf{computer} taking up about \textbf{the} size of this room, less powerful than \textbf{the} \textbf{computer} in your cell phone.
But Moore ' s Law, which is very often identified with this exponential growth, is just one example of many, because it ' s basically a property of \textbf{the} \textbf{evolutionary} process of technology.
I put 49 \textbf{famous} \textbf{computers} on this logarithmic graph — by \textbf{the} way, a straight line on a logarithmic graph is exponential growth — that ' s another exponential.
It took us three years to double our price performance of computing in 1900, two years in \textbf{the} middle ; we ' re now doubling it every one year.
And that ' s exponential growth through five different paradigms.
Moore ' s Law was just \textbf{the} last part of that, where we were shrinking transistors on an integrated circuit, but we had electro-mechanical calculators, relay-based \textbf{computers} that cracked \textbf{the} German Enigma Code, vacuum tubes in \textbf{the} 1950s predicted \textbf{the} election of Eisenhower, discreet transistors used in \textbf{the} first space flights and then Moore ' s Law.
Every time one paradigm ran out of steam, another paradigm came out of left field to continue \textbf{the} exponential growth.
They were shrinking vacuum tubes, making them smaller and smaller.
That hit a wall.
They couldn ' t shrink them and keep \textbf{the} vacuum.
Whole different paradigm — transistors came out of \textbf{the} woodwork.
In fact, when we see \textbf{the} end of \textbf{the} line for a particular paradigm, it creates research pressure to create \textbf{the} next paradigm.
And because we ' ve been predicting \textbf{the} end of Moore ' s Law for quite a long time — \textbf{the} first prediction said 2002, until now it says 2022.
But by \textbf{the} teen years, \textbf{the} features of transistors will be a few atoms in \textbf{width}, and we won ' t be able to shrink them any more.
That ' ll be \textbf{the} end of Moore ' s Law, but it won ' t be \textbf{the} end of \textbf{the} exponential growth of computing, because chips are flat.
We live in a three-dimensional world ; we might as well use \textbf{the} third \textbf{dimension}.
We will go into \textbf{the} third \textbf{dimension} and there ' s been tremendous progress, just in \textbf{the} last few years, of getting three-dimensional, self- organizing molecular circuits to work.
We ' ll have those ready well before Moore ' s Law runs out of steam.
Supercomputers — same thing.
Processor performance on Intel chips, \textbf{the} average price of a transistor — 1968, you could buy one transistor for a dollar.
You could buy 10 million in 2002.
It ' s pretty remarkable how smooth an exponential process that is.
I mean, you ' d think this is \textbf{the} result of some tabletop experiment, but this is \textbf{the} result of worldwide chaotic behavior — countries accusing each other of dumping products, IPOs, bankruptcies, marketing programs.
You would think it would be a very erratic process, and you have a very smooth outcome of this chaotic process.
Just as we can ' t predict what one molecule in a gas will do — it ' s hopeless to predict a single molecule — yet we can predict \textbf{the} properties of \textbf{the} whole gas, using thermodynamics, very accurately.
It ' s \textbf{the} same thing here.
We can ' t predict any particular project, but \textbf{the} result of this whole worldwide, chaotic, unpredictable activity of competition and \textbf{the} \textbf{evolutionary} process of technology is very predictable.
And we can predict these trends far into \textbf{the} future.
Unlike Gertrude Stein ' s roses, it ' s not \textbf{the} case that a transistor is a transistor.
As we make them smaller and less expensive, \textbf{the} electrons have less distance to travel.
They ' re faster, so you ' ve got exponential growth in \textbf{the} speed of transistors, so \textbf{the} cost of a cycle of one transistor has been coming down with a halving rate of 1. 1 years.
You add other forms of innovation and processor design, you get a doubling of price performance of computing every one year.
And that ' s basically deflation — 50 percent deflation.
And it ' s not just \textbf{computers}.
I mean, it ' s true of DNA sequencing ; it ' s true of brain scanning ; it ' s true of \textbf{the} World Wide Web.
I mean, anything that we can quantify, we have hundreds of different measurements of different, information-related measurements — capacity, adoption rates — and they basically double every 12, 13, 15 months, depending on what you ' re looking at.
In terms of price performance, that ' s a 40 to 50 percent deflation rate.
And economists have actually started worrying about that.
We had deflation during \textbf{the} Depression, but that was collapse of \textbf{the} money supply, collapse of consumer confidence, a completely different phenomena.
This is due to greater productivity, but \textbf{the} economist says,"But there ' s no way you ' re going to be able to keep up with that.
If you have 50 percent deflation, people may increase their volume 30, 40 percent, but they won ' t keep up with it."But what we ' re actually seeing is that we actually more than keep up with it.
We ' ve had 28 percent per year compounded growth in dollars in information technology over \textbf{the} last 50 years.
I mean, people didn ' t build iPods for 10, 000 dollars 10 years ago.
As \textbf{the} price performance makes new applications feasible, new applications come to \textbf{the} market.
And this is a very widespread phenomena.
Magnetic data storage — that ' s not Moore ' s Law, it ' s shrinking magnetic spots, different engineers, different companies, same exponential process.
A key revolution is that we ' re understanding our own \textbf{biology} in these information terms.
We ' re understanding \textbf{the} software programs that make our body run.
These were \textbf{evolved} in very different times — we ' d like to actually change those programs.
One little software program, called \textbf{the} \textbf{fat} insulin receptor gene, basically says,"Hold onto every calorie, because \textbf{the} next hunting season may not work out so well."That was in \textbf{the} interests of \textbf{the} species tens of thousands of years ago.
We ' d like to actually turn that program off.
They tried that in animals, and these mice ate ravenously and remained slim and got \textbf{the} health benefits of being slim.
They didn ' t get diabetes ; they didn ' t get heart disease ; they lived 20 percent longer ; they got \textbf{the} health benefits of caloric restriction without \textbf{the} restriction.
Four or five pharmaceutical companies have noticed this, felt that would be interesting drug for \textbf{the} human market, and that ' s just one of \textbf{the} 30, 000 genes that affect our biochemistry.
We were \textbf{evolved} in an era where it wasn ' t in \textbf{the} interests of people at \textbf{the} age of most people at this conference, like myself, to live much longer, because we were using up \textbf{the} precious resources which were better deployed towards \textbf{the} children and those caring for them.
So, life — long lifespans — like, that is to say, much more than 30 — weren ' t selected for, but we are learning to actually \textbf{manipulate} and change these software programs through \textbf{the} biotechnology revolution.
For example, we can inhibit genes now with RNA interference.
There are exciting new forms of gene therapy that overcome \textbf{the} problem of placing \textbf{the} genetic material in \textbf{the} right place on \textbf{the} chromosome.
There ' s actually a — for \textbf{the} first time now, something going to human trials, that actually cures pulmonary hypertension — a fatal disease — using gene therapy.
So we ' ll have not just designer babies, but designer baby boomers.
And this technology is also accelerating.
It cost 10 dollars per base pair in 1990, then a penny in 2000.
It ' s now under a 10th of a cent. \textbf{The} amount of genetic data — basically this shows that smooth exponential growth doubled every year, enabling \textbf{the} \textbf{genome} project to be completed.
Another major revolution : \textbf{the} communications revolution. \textbf{The} price performance, bandwidth, capacity of communications measured many different ways ; wired, wireless is growing exponentially. \textbf{The} Internet has been doubling in power and continues to, measured many different ways.
This is based on \textbf{the} number of hosts.
Miniaturization — we ' re shrinking \textbf{the} size of technology at an exponential rate, both wired and wireless.
These are some designs from Eric Drexler ' s book — which we ' re now showing are feasible with super-computing simulations, where actually there are scientists building molecule-scale robots.
One has one that actually walks with a surprisingly human-like gait, that ' s built out of molecules.
There are little machines doing things in experimental bases. \textbf{The} most exciting opportunity is actually to go inside \textbf{the} human body and perform therapeutic and diagnostic functions.
And this is less futuristic than it may sound.
These things have already been done in animals.
There ' s one nano-engineered \textbf{device} that cures type 1 diabetes.
It ' s \textbf{blood} cell- sized.
They put tens of thousands of these in \textbf{the} \textbf{blood} cell — they tried this in rats — it lets insulin out in a controlled fashion, and actually cures type 1 diabetes.
What you ' re watching is a design of a robotic red \textbf{blood} cell, and it does bring up \textbf{the} issue that our \textbf{biology} is actually very sub- optimal, even though it ' s remarkable in its intricacy.
Once we understand its principles of operation, and \textbf{the} pace with which we are reverse-engineering \textbf{biology} is accelerating, we can actually design these things to be thousands of times more capable.
An analysis of this respirocyte, designed by Rob Freitas, indicates if you replace 10 percent of your red \textbf{blood} cells with these robotic versions, you could do an Olympic sprint for 15 minutes without taking a breath.
You could sit at \textbf{the} bottom of your pool for four hours — so,"Honey, I ' m in \textbf{the} pool,"will take on a whole new meaning.
It will be interesting to see what we do in our Olympic trials.
Presumably we ' ll ban them, but then we ' ll have \textbf{the} specter of teenagers in their high schools gyms routinely out- performing \textbf{the} Olympic athletes.
Freitas has a design for a robotic white \textbf{blood} cell.
These are 2020-circa scenarios, but they ' re not as futuristic as it may sound.
There are four major conferences on building \textbf{blood} cell-sized \textbf{devices} ; there are many experiments in animals.
There ' s actually one going into human trial, so this is feasible technology.
If we come back to our exponential growth of computing, 1, 000 dollars of computing is now somewhere between an insect and a mouse brain.
It will intersect human intelligence in terms of capacity in \textbf{the} 2020s, but that ' ll be \textbf{the} hardware side of \textbf{the} equation.
Where will we get \textbf{the} software?
Well, it turns out we can see inside \textbf{the} human brain, and in fact not surprisingly, \textbf{the} spatial and temporal resolution of brain scanning is doubling every year.
And with \textbf{the} new generation of scanning tools, for \textbf{the} first time we can actually see individual inter-neural fibers and see them processing and signaling in real time — but then \textbf{the} question is, OK, we can get this data now, but can we understand it?
Doug Hofstadter wonders, well, maybe our intelligence just isn ' t great enough to understand our intelligence, and if we were smarter, well, then our brains would be that much more complicated, and we ' d never catch up to it.
It turns out that we can understand it.
This is a block diagram of a model and simulation of \textbf{the} human auditory cortex that actually works quite well — in applying psychoacoustic tests, gets very similar results to human auditory perception.
There ' s another simulation of \textbf{the} cerebellum — that ' s more than half \textbf{the} neurons in \textbf{the} brain — again, works very similarly to human skill formation.
This is at an early stage, but you can show with \textbf{the} exponential growth of \textbf{the} amount of information about \textbf{the} brain and \textbf{the} exponential improvement in \textbf{the} resolution of brain scanning, we will succeed in reverse-engineering \textbf{the} human brain by \textbf{the} 2020s.
We ' ve already had very good models and simulation of about 15 regions out of \textbf{the} several hundred.
All of this is driving exponentially growing economic progress.
We ' ve had productivity go from 30 dollars to 150 dollars per hour of labor in \textbf{the} last 50 years.
E-commerce has been growing exponentially.
It ' s now a trillion dollars.
You might wonder, well, wasn ' t there a boom and a bust?
That was strictly a capital-markets phenomena.
Wall Street noticed that this was a revolutionary technology, which it was, but then six months later, when it hadn ' t revolutionized all business models, they figured, well, that was wrong, and then we had this bust.
All right, this is a technology that we put together using some of \textbf{the} technologies we ' re involved in.
This will be a routine feature in a cell phone.
It would be able to translate from one language to another.
So let me just end with a couple of scenarios.
By 2010 \textbf{computers} will \textbf{disappear}.
They ' ll be so small, they ' ll be embedded in our clothing, in our environment.
Images will be written directly to our retina, providing full- immersion virtual reality, augmented real reality.
We ' ll be interacting with virtual personalities.
But if we go to 2029, we really have \textbf{the} full maturity of these trends, and you have to appreciate how many turns of \textbf{the} screw in terms of generations of technology, which are getting faster and faster, we ' ll have at that point.
I mean, we will have two-to-the-25th-power greater price performance, capacity and bandwidth of these technologies, which is pretty phenomenal.
It ' ll be millions of times more powerful than it is today.
We ' ll have completed \textbf{the} reverse-engineering of \textbf{the} human brain, 1, 000 dollars of computing will be far more powerful than \textbf{the} human brain in terms of basic raw capacity. \textbf{Computers} will combine \textbf{the} subtle pan-recognition powers of human intelligence with ways in which machines are already superior, in terms of doing analytic thinking, remembering billions of facts accurately.
Machines can share their knowledge very quickly.
But it ' s not just an alien invasion of intelligent machines.
We are going to merge with our technology.
These nano-bots I mentioned will first be used for medical and health applications : cleaning up \textbf{the} environment, providing powerful fuel cells and widely distributed decentralized solar panels and so on in \textbf{the} environment.
But they ' ll also go inside our brain, interact with our biological neurons.
We ' ve demonstrated \textbf{the} key principles of being able to do this.
So, for example, full-immersion virtual reality from within \textbf{the} nervous system, \textbf{the} nano-bots shut down \textbf{the} signals coming from your real senses, replace them with \textbf{the} signals that your brain would be receiving if you were in \textbf{the} virtual environment, and then it ' ll feel like you ' re in that virtual environment.
You can go there with other people, have any kind of experience with anyone involving all of \textbf{the} senses."Experience beamers,"I call them, will put their whole flow of sensory experiences in \textbf{the} neurological correlates of their emotions out on \textbf{the} Internet.
You can plug in and experience what it ' s like to be someone else.
But most importantly, it ' ll be a tremendous expansion of human intelligence through this direct merger with our technology, which in some sense we ' re doing already.
We routinely do intellectual feats that would be impossible without our technology.
Human life expectancy is expanding.
It was 37 in 1800, and with this sort of biotechnology, nano-technology revolutions, this will move up very rapidly in \textbf{the} years ahead.
My main message is that progress in technology is exponential, not linear.
Many — even scientists — assume a linear model, so they ' ll say,"Oh, it ' ll be hundreds of years before we have self-replicating nano-technology assembly or \textbf{artificial} intelligence."If you really look at \textbf{the} power of exponential growth, you ' ll see that these things are pretty soon at hand.
And information technology is increasingly encompassing all of our lives, from our music to our manufacturing to our \textbf{biology} to our energy to materials.
We ' ll be able to manufacture almost anything we need in \textbf{the} 2020s, from information, in very inexpensive raw materials, using nano-technology.
These are very powerful technologies.
They both empower our promise and our peril.
So we have to have \textbf{the} will to apply them to \textbf{the} right problems.
Thank you very much.

% matched lemmas: agriculture, artificial, biology, blood, computer, device, dimension, disappear, enjoy, evolution, evolutionary, evolve, explosion, famous, fat, friendly, genome, manipulate, renewable, the, theme, web, width
\end{textsample}
